\documentclass[10pt,a4paper]{amsart}
\usepackage[margin=19mm]{geometry}
\usepackage{amsmath,amssymb,mathtools}
\usepackage[scr=rsfs,cal=euler]{mathalfa}
\usepackage{xfrac}
\usepackage[unicode,hidelinks,bookmarks=false]{hyperref}
\newcommand{\ie}{i.e.\ }
\newcommand{\cf}{cf.\ }
\newcommand{\resp}{respectively\ }
\newcommand{\Subsection}{\subsection}
\providecommand{\nobreakdash}{\nobreak\hbox{-}}
\setlength{\emergencystretch}{2em}
\multlinegap0pt
\usepackage{bm}
\usepackage{bbm}
\usepackage{dsfont}
\usepackage{xspace}
\usepackage{cancel}
\usepackage{mathtools}
\usepackage{amsthm}
\makeatletter
\@ifundefined{lemma}{\newtheorem{lemma}{Lemma}}{}
\makeatother
\title[A necessary condition for the logarithmic Minkowksi problem ]{A necessary condition for the logarithmic Minkowksi problem in higher dimension}
\author{Mijia Lai, Zixiao Wang}
\date{Source: arXiv:2601.18167v1 (CC BY 4.0)}
\begin{document}
\maketitle

\noindent\emph{Source and licence.} A necessary condition for the logarithmic Minkowksi problem in higher dimension. Version arXiv:2601.18167v1 (CC BY 4.0), distributed under the Creative Commons Attribution 4.0 International licence, as recorded in the arXiv licence metadata for this exact version and shown on the abstract page https://arxiv.org/abs/2601.18167v1. Full rights notice: licenses/arxiv-2601.18167-CC-BY-4.0.txt.\par\smallskip

\noindent\textit{Editorial scope.} Counted unit: the Lemma whose source label is \texttt{None} in the article (source0.tex lines 486--491, source label \texttt{None}), delivered together with the article's own complete original proof of it (source0.tex lines 493--557, one \texttt{proof} environment of 65 lines). No mathematics is written, repaired or re-derived: statement and proof are the article's own text line for line. Required context retained uncounted: L2 (lines 326--330). Every definition, display and label of those lines is retained unchanged. Omitted from this package and not redistributed: 1--325, 331--485, 492--492, 558--562. Nothing inside a retained excerpt is omitted or altered: every retained body is delivered exactly as the article's lines are written, so no omission receipt of any kind accompanies this package. Citations: the retained excerpts carry no citation command at all, so no bibliography entry of the article is transcribed and no reference is redistributed. Reference adaptation: none was needed and none was made; every reference inside the delivered unit resolves inside it, so no unresolved reference or citation is produced. Printed slips kept literal: no printed word, symbol, hypothesis or number of the counted statement or of its proof was altered. Layout and class: the article's own class is \texttt{amsart} with packages including xy, parskip, verbatim, color, amsmath, amscd, graphicx, latexsym, hyperref, times, esint, overpic, tikz, tikz-cd; the wrapper is the lane's pinned \texttt{amsart} editorial wrapper at 10pt on A4, which restates the theorem-like environments the retained text needs and adds the article's own macro definitions that the retained text needs (none), with extra wrapper packages (bm, bbm, dsfont, xspace, cancel, mathtools). The delivered numbering is the unit's own. Assets: no retained span contains a figure environment, a graphics-inclusion command, a PDF-page inclusion, a table or a TikZ picture; no image file is copied into this package, the package declares no asset dependency and no image file is redistributed with it. Licence: version arXiv:2601.18167v1 of this article is distributed under the Creative Commons Attribution 4.0 International licence, as recorded in the arXiv licence metadata for this exact version and shown on the abstract page https://arxiv.org/abs/2601.18167v1. A full rights notice is delivered at licenses/arxiv-2601.18167-CC-BY-4.0.txt. Characters: every retained excerpt is pure ASCII, so the delivered engine, XeLaTeX with its default Latin Modern text font, reports no missing character.
\begin{lemma} \label{L2}
        \begin{align} \label{keyinequality}
\frac{\left(t^n-1\right)^{2 n}-n^2 t^{n-1}(t-1)^2\left(t^n-1\right)^{2 n-2}}{\left(t^{2 n}-n t^{n+1}+n t^{n-1}-1\right)^n} \geq 1, \quad \forall t\geq 1.
\end{align}
\end{lemma}

\begin{lemma} 
Assume $n\geq 3$, then for $t>1$, we have 
\begin{align} \notag 
\frac{\left(t^n-1\right)^{2 n}-n^2 t^{n-1}(t-1)^2\left(t^n-1\right)^{2 n-2}}{\left(t^{2 n}-n t^{n+1}+n t^{n-1}-1\right)^n} \geq 1.
\end{align}
\end{lemma}

\begin{proof}[Proof of Lemma \ref{L2}]
\ 

Set $F=\left(t^n-1\right)^{2 n}-n^2 t^{n-1}(t-1)^2\left(t^n-1\right)^{2 n-2}$ and $
G=\left(t^{2 n}-n t^{n+1}+n t^{n-1}-1\right)^n$.

We \textbf{claim} that $\frac{F}{G}$ is non increasing. Given that, notice that 
\[
\lim_{t\to \infty} \frac{F(t)}{G(t)}=1,
\]
and thus the inequality follows. 

To show that $\frac{F}{G}$ is non increasing, it suffices to show 
\[
\frac{F'}{F}\leq \frac{G'}{G}.
\]
We have
\[
\frac{F^{\prime}}{F}=n^2 \frac{t^{n-2}\left[2 t\left(t^n-1\right)^2-2 n(n-1) t^n(t-1)^2-(t-1)\left(t^n-1\right)((n+1) t-(n-1))\right]}{\left(t^n-1\right)\left[\left(t^n-1\right)^2-n^2 t^{n-1}(t-1)^2\right]},
\]
and
\[
\frac{G^{\prime}}{G}=n^2 \frac{t^{n-2}\left(2 t^{ n+1}-(n+1) t^2+(n-1) \right)}{t^{2 n}-n t^{n+1}+n t^{n-1}-1}.
\]

After simplification, it amounts to show $\forall t \geq 1$, we have 
\begin{align} \notag
(2 n-2)  &t^{3 n}+\left(2 n^3-2 n^2-2 n+2\right) t^{2 n+1}+\left(n^3-n\right) t^{2 n-2}+\left(2 n^3-2 n^2-4\right) t^{n+1} \\ \notag
& +\left(2 n^2-4 n-6\right) t^n+\left(n^3-2 n^2+n\right) t^{n-2}+(2 n+2) t \\\notag
& -\{ (2 n-2)+\left(2 n^3-2 n^2-2 n+2\right) t^{n-1}+\left(n^3-n\right) t^{n+2} \\\notag
& +\left(2 n^3-2 n^2-4\right) t^{2 n-1}+\left(2 n^2-4 n-6\right) t^{2 n}+\left(n^3-2 n^2+n\right) t^{2 n+2} \\ \notag
& +(2 n+2) t^{3 n-1} \}\geq 0.
\end{align}

Denote the polynomial on the left hand side by $p_1(t)$.  By direction computation, we find 
\[
p_1^{(m)}(1)=0, m=0,1,2.
\]
Let $p_2(t)=\frac{p_1^{(2)}(t)}{t^{n-4}}$, we again find that 
\[
p_2^{(m)}(1)=0, \quad m=0, 1, 2, 3, 4. 
\]
Let $p_3(t)=\frac{p_2^{(5)}(t)}{t^{n-5}}$, then 
\begin{align} \notag 
p_3(1)&=140n^2 (n-1)^2 (n+1)^2 (n-2)>0, \\ \notag
p_3'(1)&=140n^2 (n-1)^2 (n+1)^2 (n-2)(5n+1)>0, \\ \notag
p_3''(1)&=n^2 (n-1)^2 (n+1)^2 (n-2)(1568n^2+1344n+336)>0, \\ \notag
p_3^{(3)}(1)&=n^2 (n-1)^2 (n+1)^2 (n-2)(n+1/2)(2352n^2+1456n+1344)>0, \\ \notag
p_3^{(4)}(1)&=n^2 (n-1)^2 (n+1)^2 (n-2)(n+1/2)(2976n^3-848n^2+2608n+1056)>0, \\ \notag
p_3^{(5)}(1)&=n^2 (n-1)^2 (n+1)^2 (n+1/2)(n-1/2)(n-1/3) \\ \notag
& \cdot (n-2)(3552n^2-5742n-1728)>0.
\end{align}
Set $p_4(t)=\frac{p_3^{(5)}(t)}{t^{n-4}}$, then 
\begin{align} \notag 
p_4'(t)&\equiv (2n-2)3n(3n-1)(2n+2)(2n+1)2n\\ \notag
 &\cdot (2n-1)(2n-2)(n+2)(n+1)n(n-1)(n-2)>0.
\end{align}
Based on these computations, $p_1(t)\geq 0$, $\forall t\geq 1$. In fact, it is easy to see the inequality is strict. 

Note if $n=2$, we have the identity
\[
\frac{\left(t^n-1\right)^{2 n}-n^2 t^{n-1}(t-1)^2\left(t^n-1\right)^{2 n-2}}{\left(t^{2 n}-n t^{n+1}+n t^{n-1}-1\right)^n} \equiv  1,
\]
which explains why, in dimension two, if equality holds the convex body $K$ can generally be a trapezoid.
\end{proof}

\end{document}
